\section{Boundary kernels and the multiplier identity}\label{sec:boundary}

Our objective is a boundary formula for the nonnegative vector energy from \Cref{prop:vector}. We first construct scalar solution and interaction kernels for a general nonnegative subcritical potential on the disk. We then return to the conformal potential and derive the multiplier identity that will connect these kernels to an eigenfunction gradient.

\subsection{A subcritical disk potential}

The conformal map of \Cref{lem:conformal} defines a measure
\begin{equation}\label{eq:q}
\dd q(x)=\mu|\Phi'(x)|^2\dd x.
\end{equation}
It has a bounded positive density. A subscript $q$ on an inner product or norm denotes $L^2(q)$. Conformal invariance of the Dirichlet integral and the Dirichlet inequality on $\Omega$ give
\begin{equation}\label{eq:subcritical}
\norm\psi_{L^2(q)}^2\le d\int_\D|\grad\psi|^2,
\qquad \psi\in H^1_0(\D),\qquad d=\frac\mu{\lambda_D}<1.
\end{equation}
Until the vector-form application in \Cref{subsec:energy}, only boundedness and nonnegativity of the density and \eqref{eq:subcritical} are needed. In particular, the kernel constructions apply to any measure $q$ with those properties.
For such a general measure we fix a constant $0<d<1$ satisfying the same inequality; the zero measure also admits this choice.

The Dirichlet Green function and Poisson kernel of the disk are
\begin{align}
G(x,y)&=\frac1{2\pi}\log\left|\frac{1-x\overline y}{x-y}\right|
=\frac1{4\pi}\log\left(1+\frac{(1-|x|^2)(1-|y|^2)}{|x-y|^2}\right),\label{eq:green}\\
P_s(x)&=\frac{1-|x|^2}{2\pi|s-x|^2},\qquad s\in\Sone.\label{eq:poisson}
\end{align}
The Green function is positive and is assigned value $+\infty$ on the diagonal. Its logarithmic pole, harmonic correction, and zero boundary data give the Green identity. Also $\int_{\Sone}P_s(x)\dd s=1$.

\begin{lemma}\label{lem:green-operator}
The operator
\[
Tf(x)=\int_\D G(x,y)f(y)\dd q(y)
\]
is a bounded self-adjoint operator on $L^2(q)$ with $\norm T\le d<1$. The same statement holds with $q$ replaced by any measure $0\le\rho\le q$.
\end{lemma}
\begin{proof}
Write $\dd q=a\dd x$, with $a$ bounded. For $f\in L^2(q)$, the source $af$ lies in $L^2(\dd x)$, since $\int|af|^2\le\norm a_\infty\int|f|^2\dd q$. Its Dirichlet solution $\psi=Tf$ satisfies
\[
\int_\D|\grad\psi|^2=\ip{\psi}{f}_q.
\]
Combining this with \eqref{eq:subcritical} gives
\[
\norm\psi_q^2\le d\norm\psi_q\norm f_q,
\]
hence the norm bound. Symmetry follows from the Green kernel, or by testing the two Poisson equations against each other. The Green formula can first be proved for smooth compactly supported sources and then passed to $L^2(\dd x)$ by approximation; $G(x,\cdot)\in L^2(\dd x)$ gives pointwise evaluation. These sections vary continuously in $L^2$ on compact interior sets, so this is the continuous interior representative. For $\rho\le q$, its density remains bounded and its weighted Dirichlet inequality follows from \eqref{eq:subcritical}, proving the same result.
\end{proof}

For Lipschitz scalar or vector boundary data $h$, let
\[
Ah(x)=\int_{\Sone}P_s(x)h(s)\dd s.
\]
This harmonic extension belongs to $H^1(\D)$. The unique weak solution with trace $h$ of
\[
-\Delta W_h\dd x=W_h\dd q
\]
is
\begin{equation}\label{eq:extension}
W_h=Ah+T(I-T)^{-1}Ah.
\end{equation}
Indeed the correction has zero trace and solves the required Poisson equation. Uniqueness follows from \eqref{eq:subcritical}. The construction is componentwise for vectors.

\subsection{Positive kernels and their estimates}

For $p\in\D$ and distinct $s,t\in\Sone$, define
\begin{align}
K(p,s)&=P_s(p)+\sum_{n\ge1}(T^nP_s)(p),\label{eq:K}\\
R(s,t)&=\sum_{n\ge0}\int_\D P_s\,T^nP_t\dd q,\label{eq:R}\\
N(s,t)&=N_0(s,t)+R(s,t),\qquad N_0(s,t)=\frac1{\pi|s-t|^2}.\label{eq:N}
\end{align}
Powers applied to $P_s$ mean iterated nonnegative kernel integrals, not an a priori application of an $L^2$ operator to $P_s$.

\begin{lemma}\label{lem:kernel-estimates}
The kernels above are finite at the stated points. For fixed $p$, $K(p,\cdot)$ is positive and bounded. The kernel $R$ is symmetric, nonnegative, and integrable on $\Sone\times\Sone$, with
\begin{equation}\label{eq:R-integral}
\iint R(s,t)\dd s\dd t\le\frac{q(\D)}{1-d}.
\end{equation}
For Lipschitz boundary data,
\begin{align}
W_h(p)&=\int_{\Sone}K(p,s)h(s)\dd s,\label{eq:K-representation}\\
\ip{Ah}{(I-T)^{-1}Ak}_q
&=\iint R(s,t)h(s)\cdot k(t)\dd s\dd t.\label{eq:R-representation}
\end{align}
\end{lemma}
\begin{proof}
The elementary estimates
\[
P_s(y)\le\frac C{|y-s|},\qquad
0\le G(x,y)\le C\log\frac2{|x-y|}
\]
give uniform bounds for $P_s$ in $L^{3/2}(\dd x)$ and for $G(x,\cdot)$ in $L^3(\dd x)$. H\"older and the bounded density imply
\[
\sup_{s\in\Sone}\norm{TP_s}_{L^\infty(\D)}<\infty.
\]
The section $P_s$ need not belong to $L^2(q)$, but one Green integration gives $F_s=TP_s\in L^\infty(\D)\subset L^2(q)$. We can now apply the $L^2(q)$ contraction estimate to the remaining iterates. For $n\ge2$,
\begin{equation}\label{eq:iterate-bound}
|(T^nP_s)(p)|
=|\ip{G(p,\cdot)}{T^{n-2}F_s}_q|
\le\norm{G(p,\cdot)}_q\,d^{n-2}\norm{F_s}_q
\le C d^{n-2}.
\end{equation}
The constants can be chosen uniformly in $p$ and $s$. This proves the assertion about $K$.

For distinct $s,t$, the product $P_sP_t$ is integrable: near either boundary pole, the other factor is bounded, and $|y-s|^{-1}$ is locally integrable in two dimensions. The terms with $n\ge1$ in \eqref{eq:R} are summable by the first smoothing estimate and \eqref{eq:iterate-bound}, since $\int P_s\dd q$ is uniformly bounded. Reversing the nonnegative chains proves symmetry. Tonelli's theorem and $\int P_s\dd s=1$ give
\[
\iint R(s,t)\dd s\dd t
=\sum_{n\ge0}\ip\one{T^n\one}_q
\le\sum_{n\ge0}d^n\norm\one_q^2,
\]
which is \eqref{eq:R-integral}. Diagonal values of $R$ play no role in these boundary integrals. Expanding the Neumann series in \eqref{eq:extension} and using the positive kernels gives \eqref{eq:K-representation} and \eqref{eq:R-representation}; the same bounds justify the exchanges for signed data by absolute domination.
\end{proof}

\subsection{The boundary energy}\label{subsec:energy}

Return now to $q$ in \eqref{eq:q}. A vector boundary function $h$ is called tangent if $h(s)$ is tangent to $\partial\Omega$ at $\Phi(s)$. Extend its fixed Cartesian components by \eqref{eq:extension}, and set
\[
X_h=W_h\circ\Phi^{-1},\qquad
E(h)=\Q(X_h)-\mu\norm{X_h}_{L^2(\Omega)}^2.
\]
We write $E(h,k)$ for its bilinear form. \Cref{prop:vector} gives $E\ge0$ on tangent Lipschitz boundary data, and
\begin{equation}\label{eq:E-nullspace}
E(h)=0\quad\Longleftrightarrow\quad
X_h=\grad v\text{ for some }v\in\V.
\end{equation}

\begin{lemma}\label{lem:boundary-energy}
For tangent Lipschitz data $h,k$,
\begin{align}
E(h,k)={}&\frac12\iint N_0(s,t)
(h(s)-h(t))\cdot(k(s)-k(t))\dd s\dd t\notag\\
&-\iint R(s,t)h(s)\cdot k(t)\dd s\dd t
+\int_{\Sone}|\Phi'|(\kappa\circ\Phi)h\cdot k\dd s.
\label{eq:E-boundary}
\end{align}
\end{lemma}
\begin{proof}
By \Cref{lem:curvature} and conformal invariance of the scalar Dirichlet integral applied to each Cartesian component, the bulk terms of $E$ are
\[
\int_\D\grad W_h:\grad W_k-\ip{W_h}{W_k}_q.
\]
The equation for $W_k$ permits replacing $W_h$ here by $Ah$, since their difference has zero trace. Harmonicity of $Ah$ then replaces $W_k$ by $Ak$ in the gradient term. Since $W_k=(I-T)^{-1}Ak$ in $L^2(q)$, this gives
\[
\int_\D\grad Ah:\grad Ak-\ip{Ah}{(I-T)^{-1}Ak}_q.
\]
The boundary curvature term changes by the factor $|\Phi'|$, and \eqref{eq:R-representation} identifies the second term.

For the harmonic energy,
\begin{equation}\label{eq:harmonic-energy}
\int_\D\grad Ah:\grad Ak
=\frac12\iint N_0(s,t)(h(s)-h(t))\cdot(k(s)-k(t))\dd s\dd t.
\end{equation}
Indeed both sides diagonalize in Fourier modes. The Hermitian energy of $s^m$ is $2\pi|m|$, while
\[
\iint\frac{|s^m-t^m|^2}{|s-t|^2}\dd s\dd t=(2\pi)^2|m|
\]
by a geometric sum, verifying the normalization. Smooth convolution approximation, with uniformly bounded Lipschitz constants and convergence in $H^{1/2}$, extends the identity to the stated data. This proves \eqref{eq:E-boundary}.
\end{proof}

Fix $0\ne u\in\V$ and let
\begin{equation}\label{eq:g}
g(s)=(\grad u)(\Phi(s)),\qquad s\in\Sone.
\end{equation}
It is tangent by the Neumann condition. Differentiation of the Euclidean Helmholtz equation commutes with its Cartesian derivatives, so uniqueness gives
\begin{equation}\label{eq:gradient-extension}
W_g=(\grad u)\circ\Phi,
\qquad E(g)=0.
\end{equation}

\begin{proposition}[Multiplier identity]\label{prop:multiplier}
For every real Lipschitz function $b$ on $\Sone$,
\begin{equation}\label{eq:multiplier}
E(bg)=\frac12\iint N(s,t)(b(s)-b(t))^2
g(s)\cdot g(t)\dd s\dd t.
\end{equation}
\end{proposition}
\begin{proof}
Since $E$ is positive semidefinite and $E(g)=0$, its bilinear form annihilates $g$: this follows by considering $E(g+th)$ for real $t$. In particular, $E(g,b^2g)=0$. Subtract \eqref{eq:E-boundary} for this pair from the formula for $E(bg,bg)$. The curvature terms cancel. The harmonic term gives
\[
\frac12N_0(s,t)(b(s)-b(t))^2g(s)\cdot g(t)
\]
in the double integral. Symmetry of $R$ gives the same expression with $R$ in place of $N_0$. Their sum is \eqref{eq:multiplier}. All terms are integrable: the squared Lipschitz difference cancels the singularity of $N_0$, and $R$ is integrable.
\end{proof}
